\documentclass[11pt,a4paper]{article}

\usepackage[utf8]{inputenc}
\usepackage[T1]{fontenc}
\usepackage{amsmath,amssymb,amsthm}
\usepackage{geometry}
\usepackage{listings}
\usepackage{xcolor}
\usepackage{hyperref}
\usepackage{booktabs}
\usepackage{float}
\usepackage{enumitem}

\geometry{margin=2.5cm}

\hypersetup{
    colorlinks=true,
    linkcolor=blue!60!black,
    urlcolor=blue!60!black,
    citecolor=blue!60!black
}

\definecolor{codebg}{rgb}{0.96,0.96,0.94}
\definecolor{codekw}{rgb}{0.13,0.29,0.53}
\definecolor{codestr}{rgb}{0.58,0.13,0.13}
\definecolor{codecom}{rgb}{0.35,0.45,0.35}

\lstdefinestyle{matlab}{
    language=Matlab,
    backgroundcolor=\color{codebg},
    basicstyle=\ttfamily\small,
    keywordstyle=\color{codekw}\bfseries,
    stringstyle=\color{codestr},
    commentstyle=\color{codecom}\itshape,
    numbers=left,
    numberstyle=\tiny\color{gray},
    numbersep=8pt,
    showstringspaces=false,
    breaklines=true,
    frame=single,
    rulecolor=\color{gray!40},
    columns=fullflexible,
    keepspaces=true,
    upquote=true
}

% Tables of this document are numbered A1, A2, ... so that they cannot be
% confused with Tables 1 and 2 of Shu (1977).
\renewcommand{\thetable}{A\arabic{table}}

% Placeholder for information that only the author can supply.
\newcommand{\todo}[1]{\textcolor{red}{[TODO: #1]}}

\title{\bfseries Mathematical Explanation of the Spectral Solver\\
for Shu (1977) Self-Similar Collapse}
\author{--\\[2pt] Numerical cross-validation project}
\date{\today}

\begin{document}

\maketitle

\begin{abstract}
This document explains the mathematical content of each source file of an
independent Chebyshev spectral solver for the self-similar collapse
equations of Shu (1977, \emph{ApJ} \textbf{214}, 488). The code is written
in GNU Octave and reproduces Shu's Tables~1 and~2. The purpose of this
document is to make the numerical implementation traceable from the
published equations to the corresponding source-code operations. The
Runge--Kutta integration used by Shu is replaced here by a Chebyshev
collocation formulation followed by a nonlinear Newton solve. The
resulting solution reproduces the tabulated values of Shu with a maximum
relative difference of 0.293\% for Table~1 and 0.618\% for Table~2
(Sec.~\ref{sec:validation}).
\end{abstract}

\tableofcontents
\newpage

%=======================================================================
\section{Governing equations}
%=======================================================================

Shu (1977) writes the similarity ansatz
\begin{equation}
\rho(r,t) = \frac{\alpha(x)}{4\pi G t^{2}}, \qquad
u(r,t)    = a\,v(x), \qquad
M(r,t)    = \frac{a^{3}t}{G}\,m(x), \qquad
x = \frac{r}{a t},
\label{eq:ansatz}
\end{equation}
where $a$ is the isothermal sound speed, $G$ the gravitational constant,
$\rho$ the mass density, $u$ the radial velocity (negative for inflow),
$M$ the mass interior to radius $r$, and $t$ the physical time.
Substituting into the continuity and Euler equations yields the two
coupled ordinary differential equations (Shu 1977, eqs.~11--12):
\begin{equation}
\boxed{
\bigl[(x-v)^{2}-1\bigr]\frac{dv}{dx}
=
\left[\alpha(x-v)-\frac{2}{x}\right](x-v)
}
\label{eq:odev}
\end{equation}

\begin{equation}
\boxed{
\bigl[(x-v)^{2}-1\bigr]
\frac{1}{\alpha}\frac{d\alpha}{dx}
=
\left[\alpha-\frac{2(x-v)}{x}\right](x-v)
}
\label{eq:odea}
\end{equation}
together with the algebraic mass relation (Shu 1977, eq.~10)
\begin{equation}
m(x)=x^{2}\alpha(x)\bigl(x-v(x)\bigr).
\label{eq:mass}
\end{equation}
The coefficient multiplying the derivatives vanishes on the
\emph{critical (sonic) curve}
\begin{equation}
(x-v)^{2}=1
\qquad\Longleftrightarrow\qquad
v=x-1
\quad(\text{for } x-v=1).
\label{eq:critcurve}
\end{equation}
A physical solution must cross this curve smoothly. This requires the
right-hand sides of \eqref{eq:odev}--\eqref{eq:odea} to vanish at the
crossing point as well, so the singularity of the equations is removable
for the physical solution.

%=======================================================================
\section{Numerical formulation}
%=======================================================================

The numerical procedure of Shu (1977) starts from asymptotic conditions
and integrates the ordinary differential equations with a four-step
Runge--Kutta method. The present implementation instead follows
\begin{equation}
\boxed{
\text{Shu equations}
\;\longrightarrow\;
\text{spectral collocation}
\;\longrightarrow\;
\text{nonlinear residual equations}
\;\longrightarrow\;
\text{Newton solve}.
}
\label{eq:pipeline}
\end{equation}

It does not reproduce the Runge--Kutta algorithm line by line. It solves
the same mathematical equations with a different discretization and a
different nonlinear solution method. This difference is what makes the
comparison a meaningful cross-check: the two sets of numerical solutions
were obtained independently. Two changes of variables
(Sec.~\ref{sec:logvars}) make the nonlinear boundary-value problem more
suitable for spectral collocation.

%=======================================================================
\section{Logarithmic change of variables}\label{sec:logvars}
%=======================================================================

\subsection{Logarithmic coordinate}

The similarity variable $x$ spans several decades
($x_{\min}\sim10^{-3}$ and $x_{\max}\sim20$ in the collapse calculation).
Because $t$ denotes physical time in Eq.~\eqref{eq:ansatz}, the
logarithmic coordinate is called $\tau$ here (it is called \texttt{t} in
the code, hence \texttt{t\_min}, \texttt{t\_max} and \texttt{Dt}):
\begin{equation}
\tau=\ln x,
\qquad
\frac{d}{dx}=\frac{1}{x}\frac{d}{d\tau}.
\label{eq:logcoord}
\end{equation}
The computational domain is the finite interval
\[
[\tau_{\min},\tau_{\max}]
=
[\ln x_{\min},\ln x_{\max}],
\]
which is mapped linearly onto the Chebyshev interval $[-1,1]$. A grid in
$\tau$ spaces its nodes by equal ratios in $x$, so each decade of $x$ is
resolved comparably; a grid uniform in $x$ would place almost all nodes at
large~$x$. (The Chebyshev nodes in $\tau$ additionally cluster near both
ends of the interval.)

\subsection{Log-density variable}

The reduced density has the asymptotic behavior
\begin{equation}
\alpha(x)
\sim
\sqrt{\frac{m_0}{2x^3}}
\qquad (x\rightarrow0),
\label{eq:asym0}
\end{equation}
and
\begin{equation}
\alpha(x)\sim\frac{A}{x^2}
\qquad (x\rightarrow\infty).
\label{eq:asyminf}
\end{equation}
Here $m_0=\lim_{x\to0}m(x)$ is the dimensionless mass-accretion
parameter, so that $dM/dt=m_0a^3/G$ at small radii, and $A$ is the
coefficient of the large-$x$ asymptote. Since $\alpha$ becomes large
toward the small-$x$ end of the domain, the code uses
\begin{equation}
w(x)=\ln\alpha(x),
\qquad
\alpha=e^w.
\label{eq:logdensity}
\end{equation}
This has two numerical advantages:
\begin{enumerate}
\item positivity of $\alpha$ is automatic because $e^w>0$;
\item the dynamic range of the unknown density is reduced.
\end{enumerate}
Using Eqs.~\eqref{eq:logcoord} and~\eqref{eq:logdensity}, the original
ODE system becomes
\begin{equation}
\boxed{
\bigl[(x-v)^2-1\bigr]\frac{dv}{d\tau}
=
\bigl[\alpha x(x-v)-2\bigr](x-v)
}
\label{eq:odevt}
\end{equation}
and
\begin{equation}
\boxed{
\bigl[(x-v)^2-1\bigr]\frac{dw}{d\tau}
=
\bigl[\alpha x-2(x-v)\bigr](x-v),
}
\label{eq:odewt}
\end{equation}
where
\begin{equation}
x=e^\tau,
\qquad
\alpha=e^w.
\end{equation}
These are the equations enforced by the spectral residual functions.

%=======================================================================
\section{The Chebyshev differentiation matrix --- \texttt{sc\_cheb.m}}
%=======================================================================
The file \texttt{sc\_cheb.m} constructs the Chebyshev differentiation
matrix. (The variable \texttt{x} in this listing holds the nodes $\xi_j$,
not the similarity variable.)
\begin{lstlisting}[style=matlab]
function [D, x] = sc_cheb(N)
    x = cos(pi * (0:N)' / N);
    c = [2; ones(N-1,1); 2] .* (-1).^(0:N)';
    X = repmat(x, 1, N+1);
    dX = X - X';
    D  = (c * (1./c)') ./ (dX + eye(N+1));
    D  = D - diag(sum(D, 2));
end
\end{lstlisting}
The nodes are the Chebyshev--Gauss--Lobatto nodes
\begin{equation}
\xi_j=\cos\left(\frac{\pi j}{N}\right),
\qquad
j=0,\ldots,N.
\label{eq:cgl}
\end{equation}
They cluster near $\xi=\pm1$, which is useful when high accuracy is
required near the endpoints of the computational interval. For a vector
of nodal values $\mathbf f$, the differentiation matrix satisfies
\begin{equation}
(D\mathbf f)_i
\approx
\left.\frac{df}{d\xi}\right|_{\xi_i}.
\label{eq:Dxi}
\end{equation}
The implementation uses the standard barycentric-weight construction.
The final diagonal correction makes the derivative of a constant function
numerically zero.

\subsection{Mapping between $\xi$ and $\tau$}

The computational coordinate $\tau$ is related linearly to
$\xi\in[-1,1]$:
\begin{equation}
\tau=
\frac{\tau_{\max}-\tau_{\min}}{2}\,\xi
+
\frac{\tau_{\max}+\tau_{\min}}{2}.
\end{equation}
Therefore,
\begin{equation}
\frac{d}{d\tau}
=
\frac{2}{\tau_{\max}-\tau_{\min}}\frac{d}{d\xi}.
\label{eq:dt}
\end{equation}
The code constructs
\begin{lstlisting}[style=matlab]
Dt = D * 2 / (t_max - t_min);
\end{lstlisting}
so that
\begin{equation}
D_\tau\mathbf f
\approx
\frac{d\mathbf f}{d\tau}.
\label{eq:Dt}
\end{equation}
The nodes are ordered from $\xi_0=+1$ (that is, $\tau_{\max}$ and
$x_{\max}$) to $\xi_N=-1$ ($\tau_{\min}$ and $x_{\min}$). Node $j=0$
is therefore the outer boundary.

%=======================================================================
\section{The residual equations}\label{sec:residuals}
%=======================================================================

The spectral method replaces the differential equations with a
finite-dimensional nonlinear algebraic system. Bold symbols denote
vectors of nodal values; plain $v$, $w$ denote the functions. The unknown continuous profiles are discretized by their nodal values
\begin{equation}
\mathbf v=(v_0,\ldots,v_N)^T,
\qquad
\mathbf w=(w_0,\ldots,w_N)^T,
\end{equation}
where $w_j=\ln\alpha_j$. These are concatenated into a single state vector of $2(N+1)$ nodal unknowns
\begin{equation}
\mathbf u=
\begin{bmatrix}
\mathbf v\\
\mathbf w
\end{bmatrix}.
\label{eq:statevector}
\end{equation}
These are nodal values, not expansion coefficients of a Chebyshev
polynomial. The differentiation matrix provides the derivatives of the
interpolating polynomial at the collocation nodes.

Two features of the formulation are essential. First, the residuals are
enforced at every node except the outer one, including nodes near the
critical curve~\eqref{eq:critcurve} where the coefficient of the
derivative vanishes; this forces the interpolating polynomial to satisfy
the regularity condition there, which selects the physical solution.
Second, no boundary condition is imposed at $x_{\min}$. The behavior at
small $x$ is a result of the calculation, and it is compared with the
free-fall behavior
\begin{equation}
\alpha\propto x^{-3/2},
\qquad
v\propto-x^{-1/2}
\qquad
(x\rightarrow0).
\label{eq:freefallbehavior}
\end{equation}
Indices: the nodes are numbered $j=0,\ldots,N$ in the text, while the
code uses $1,\ldots,N+1$, so the outer node $j=0$ is element \texttt{(1)}
in the code.

%=======================================================================
\subsection{Collapse problem --- \texttt{sc\_res\_collapse.m}}
%=======================================================================
The residual function first computes the spectral derivatives. Here
\texttt{x} is the vector of nodal similarity coordinates $x=e^\tau$,
\texttt{alpha} is the vector $e^{w}$, and \texttt{Dt} is the matrix of
Eq.~\eqref{eq:Dt}:
\begin{lstlisting}[style=matlab]
dv   = Dt * v;
dw   = Dt * w;
xv   = x - v;
coef = xv.^2 - 1;

R1 = coef .* dv - (alpha .* x .* xv - 2) .* xv;
R2 = coef .* dw - (alpha .* x        - 2 .* xv) .* xv;
\end{lstlisting}
These expressions correspond directly to
\begin{align}
R_1 &=
\bigl[(x-v)^2-1\bigr]\frac{dv}{d\tau}
-
\bigl[\alpha x(x-v)-2\bigr](x-v),
\label{eq:R1}\\
R_2 &=
\bigl[(x-v)^2-1\bigr]\frac{dw}{d\tau}
-
\bigl[\alpha x-2(x-v)\bigr](x-v).
\label{eq:R2}
\end{align}
At every node other than the outer one, the numerical solution must
satisfy
\begin{equation}
R_1=0,
\qquad
R_2=0.
\end{equation}

\subsubsection{Outer asymptotic boundary condition}

At the outer node $j=0$ ($x=x_{\max}$), the code replaces the two ODE
residuals by the asymptotic boundary conditions, Eq.~(19) of
Shu (1977):
\begin{lstlisting}[style=matlab]
alpha_inf = A/x_max^2 - A*(A-2)/(2*x_max^4);
v_inf     = -(A-2)/x_max - (1-A/6)*(A-2)/x_max^3;

R1(1) = v(1) - v_inf;
R2(1) = w(1) - log(alpha_inf);
\end{lstlisting}
Accordingly,
\begin{align}
R_1\big|_{j=0}
&=
v_0
-
\left[
-\frac{A-2}{x_{\max}}
-
\frac{(1-A/6)(A-2)}{x_{\max}^3}
\right],
\label{eq:bcv}\\
R_2\big|_{j=0}
&=
w_0
-
\ln\left[
\frac{A}{x_{\max}^2}
-
\frac{A(A-2)}{2x_{\max}^4}
\right].
\label{eq:bca}
\end{align}
These equations place the asymptotic solution at the truncated outer
boundary.

%=======================================================================
\subsection{Expansion-wave problem --- \texttt{sc\_res\_expwave.m}}
%=======================================================================

For $A=2$ the outer asymptote~\eqref{eq:asyminf} reduces to the static
singular isothermal sphere $\alpha=2/x^2$, $v=0$. The expansion-wave
solution is the limiting case $A\rightarrow2^+$ of the collapse family.
Its outer boundary is the critical point $x=1$, and the solution is
computed on $x_{\min}\le x\le1$ with $x_{\min}=10^{-3}$. The boundary
conditions at the outer node are
\begin{lstlisting}[style=matlab]
R1(1) = v(1);
R2(1) = w(1) - log(2);
\end{lstlisting}
that is,
\begin{equation}
v(x{=}1)=0,
\qquad
\alpha(x{=}1)=2.
\label{eq:criticalBC}
\end{equation}
These are the critical-point values of the expansion-wave solution. With
them, Eq.~\eqref{eq:mass} gives $m(x{=}1)=2$. As in the collapse
problem, no condition is imposed at $x_{\min}$.

%=======================================================================
\section{The Newton solver --- \texttt{sc\_newton\_cs.m}}
%=======================================================================

After collocation, the differential equations and boundary conditions
form a nonlinear algebraic system
\begin{equation}
\mathbf R(\mathbf u)=\mathbf 0,
\end{equation}
where $\mathbf u$ is defined in Eq.~\eqref{eq:statevector}. The file
\texttt{sc\_newton\_cs.m} solves this system.

%=======================================================================
\subsection{Complex-step Jacobian}
%=======================================================================

The Jacobian
\begin{equation}
J=\frac{\partial\mathbf R}{\partial\mathbf u}
\end{equation}
is evaluated with the complex-step method. For a residual function that
is analytic in its arguments,
\begin{equation}
\frac{\partial R_i}{\partial u_j}
\approx
\frac{
\operatorname{Im}
\left[
R_i(\mathbf u+i h\,\mathbf e_j)
\right]
}{h},
\label{eq:complexstep}
\end{equation}
where $\mathbf e_j$ is the $j$-th unit vector and $h$ is a small real
step ($h=10^{-30}$ in the code). The implementation uses
\begin{lstlisting}[style=matlab]
for j = 1:n
    up     = u;
    up(j)  = up(j) + 1i * h;
    J(:,j) = imag(resfun(up)) / h;
end
\end{lstlisting}
The complex-step method involves no subtraction of nearly equal numbers,
so it avoids the cancellation error that limits ordinary finite
differences and yields Jacobians accurate to near machine precision. The
Jacobian has dimension
\begin{equation}
2(N+1)\times2(N+1).
\end{equation}

%=======================================================================
\subsection{Row scaling}
%=======================================================================

The residual vector contains equations with different numerical scales.
The implementation scales each residual row by the largest-magnitude
entry in the corresponding Jacobian row:
\begin{lstlisting}[style=matlab]
rs = max(abs(J),[],2);
rs(rs < 1e-30) = 1;

Js = J ./ rs;
Rs = R ./ rs;
\end{lstlisting}
If $r_i$ denotes the scaling factor for row $i$, the scaled Newton
system is
\begin{equation}
\operatorname{diag}(1/r_i)\,J\,\delta\mathbf u
=
-\operatorname{diag}(1/r_i)\,\mathbf R.
\label{eq:scaled}
\end{equation}
The same non-singular row scaling is applied to both sides, so the
solution of this linear system is unchanged. The scaling does, however,
improve the balance between equations of different magnitudes. (Because
of the Tikhonov term below, the step actually computed is not exactly
invariant under the scaling; with $\epsilon=10^{-10}$ the effect is
small.)

%=======================================================================
\subsection{Tikhonov-regularized least-squares Newton step}
%=======================================================================

The Newton correction is obtained from the regularized least-squares
problem
\begin{equation}
\min_{\delta\mathbf u}
\left\|
J_s\,\delta\mathbf u+\mathbf R_s
\right\|_2^2
+
\epsilon\,\|\delta\mathbf u\|_2^2,
\label{eq:tikhonov}
\end{equation}
which is equivalent to the augmented system
\begin{equation}
\begin{bmatrix}
J_s\\
\sqrt{\epsilon}\,I_n
\end{bmatrix}
\delta\mathbf u
=
-
\begin{bmatrix}
\mathbf R_s\\
\mathbf 0
\end{bmatrix},
\label{eq:augmented}
\end{equation}
where $I_n$ is the $n\times n$ identity matrix with $n=2(N+1)$. The code
solves this system directly ($\epsilon$ is \texttt{eps\_reg}, not
Octave's \texttt{eps}):
\begin{lstlisting}[style=matlab]
du = -[Js; sqrt(eps_reg)*eye(n)] ...
       \ [Rs; zeros(n,1)];
\end{lstlisting}
This avoids forming the normal equations
$J_s^{T}J_s+\epsilon I_n$, which would generally worsen conditioning.
The regularization parameter is $\epsilon=10^{-10}$. It does not change
the target equations; it only stabilizes the linear solve when the
Jacobian is poorly conditioned, as can happen near the critical curve.
When the Jacobian is well conditioned, the step approaches the ordinary
Newton correction.

%=======================================================================
\subsection{Backtracking line search}
%=======================================================================

The Newton correction is damped by a backtracking line search. Starting
from $\lambda=1$, the code halves the step (at most 30 times) until the
residual satisfies
\begin{lstlisting}[style=matlab]
lam = 1;
for k = 1:30
    if norm(resfun(u + lam*du)) ...
            < nR*(1 - 1e-4*lam)
        break;
    end
    lam = lam/2;
end

u = u + lam*du;
\end{lstlisting}
The acceptance condition is
\begin{equation}
\|\mathbf R(\mathbf u+\lambda\,\delta\mathbf u)\|
<
\|\mathbf R(\mathbf u)\|\left(1-c\lambda\right),
\qquad
c=10^{-4}.
\label{eq:armijo}
\end{equation}
The line search prevents a large Newton step from leaving the basin of
convergence.

%=======================================================================
\subsection{Convergence criterion}
%=======================================================================

The nonlinear iteration is stopped when
\begin{equation}
\frac{\|\mathbf R(\mathbf u^k)\|}{\|\mathbf R(\mathbf u^0)\|}
<
\mathrm{tol},
\qquad
\mathrm{tol}=10^{-10}.
\label{eq:tol}
\end{equation}
This criterion measures only the algebraic convergence of the Newton
iteration, namely the reduction of the residual relative to its initial
value. It says nothing about the discretization error of the collocation
scheme or about the error caused by truncating the domain at $x_{\min}$
and $x_{\max}$. A small residual therefore does not establish agreement
with Shu's tabulated values; that comparison is made separately in
Sec.~\ref{sec:validation}.

%=======================================================================
\section{Grid refinement and continuation}
%=======================================================================

The solver refines the grid in the polynomial degree $N$ and uses
continuation in the parameter $A$.

%=======================================================================
\subsection{Grid refinement in $N$}\label{sec:contN}
%=======================================================================

The polynomial degree is increased through
\begin{lstlisting}[style=matlab]
Ns = [30, 60, 120];
\end{lstlisting}
At each stage, the solution from the previous grid is interpolated to
the new grid in the coordinate $\tau$ and used as the initial guess:
\begin{equation}
\mathbf u_0^{(k+1)}
=
\mathcal P_{N_k\rightarrow N_{k+1}}
\left[\mathbf u^{(k)}\right].
\label{eq:Ncontinuation}
\end{equation}
This avoids restarting the nonlinear solve from a crude initial guess at
every polynomial degree. Because the loop over $N$ is embedded inside
\texttt{sc\_solve\_collapse}, the sequence $N=30\to60\to120$ is repeated
for every value of $A$. Each Newton solve typically requires only a
small number of iterations at each stage, since the initial guess is
already close to the converged solution.

%=======================================================================
\subsection{Continuation in $A$}\label{sec:contA}
%=======================================================================

For the collapse solutions, the solver proceeds through
\begin{lstlisting}[style=matlab]
A_list = [2.2 2.4 2.6 2.8 3.0 3.2 3.4 3.6 3.8 4.0];
\end{lstlisting}
The solution for one value of $A$ initializes the next. Here the code
interpolates in $x$ with cubic splines:
\begin{lstlisting}[style=matlab]
v0 = interp1(x_prev, v_prev,     x, 'spline');
a0 = interp1(x_prev, alpha_prev, x, 'spline');
u  = [v0(:); log(a0(:))];
\end{lstlisting}
(\texttt{a0} is the interpolated density $\alpha$, not the sound speed
$a$.) This keeps successive nonlinear problems close to one another and
makes the solutions for larger $A$ easier to obtain.

%=======================================================================
\subsection{Initial guess}
%=======================================================================

For the first collapse calculation and for the expansion-wave
calculation, the solver constructs an initial approximation from
analytic limiting forms. The free-fall approximation is
\begin{equation}
\alpha_{\mathrm{ff}}(x)
=
\sqrt{\frac{m_0}{2x^3}},
\qquad
v_{\mathrm{ff}}(x)
=
-\sqrt{\frac{2m_0}{x}},
\label{eq:freefall}
\end{equation}
and the large-$x$ asymptotic approximation is
\begin{equation}
\alpha_{\infty}(x)
=
\frac{A}{x^2},
\qquad
v_{\infty}(x)
=
-\frac{A-2}{x}.
\label{eq:outerinitial}
\end{equation}
For the collapse calculation, the two forms are blended with the weight
\begin{equation}
\beta(x)=\frac{x^2}{1+x^2},
\end{equation}
where the density is blended in log space:
\begin{align}
\ln\alpha_0
&=
(1-\beta)\ln\alpha_{\mathrm{ff}}
+
\beta\ln\alpha_{\infty},
\\
v_0
&=
(1-\beta)v_{\mathrm{ff}}
+
\beta v_{\infty}.
\end{align}
For the expansion-wave calculation, the outer form is the hydrostatic
solution
\begin{equation}
\alpha_{\mathrm{st}}=\frac{2}{x^2},
\end{equation}
and the densities are blended linearly in $\alpha$ with the weight
$w_b = x/x_{\max}$ (with $x_{\max}=1$), so that
\begin{equation}
\alpha_0 = (1-w_b)\,\alpha_{\mathrm{ff}} + w_b\,\alpha_{\mathrm{st}},
\qquad
v_0 = (1-w_b)\,v_{\mathrm{ff}};
\end{equation}
the logarithm is then taken to form the initial $w$.

%=======================================================================
\section{Equation-to-code traceability}
%=======================================================================

Table~\ref{tab:traceability} maps each mathematical object to the code
that implements it, with equation numbers.

\begin{table}[H]
\centering
\renewcommand{\arraystretch}{1.35}
\begin{tabular}{@{}p{4.6cm}p{2.2cm}p{7.4cm}@{}}
\toprule
\textbf{Mathematical object} & \textbf{Equation} & \textbf{Implementation} \\
\midrule

Similarity variable $x=r/(at)$
& \eqref{eq:ansatz}
&
Used throughout the solver; the numerical coordinate is
$\tau=\ln x$.
\\

Logarithmic coordinate $\tau=\ln x$
& \eqref{eq:logcoord}
&
Constructed in the solver routines; the coordinate on which the
differentiation matrix acts.
\\

Log-density $w=\ln\alpha$
& \eqref{eq:logdensity}
&
The state vector stores $w$ rather than $\alpha$; $\alpha=e^w$ is
reconstructed inside the residual functions.
\\

Chebyshev--Gauss--Lobatto nodes
& \eqref{eq:cgl}
&
\texttt{sc\_cheb.m}
\\

Derivative $d/d\xi$
& \eqref{eq:Dxi}
&
Matrix \texttt{D} returned by \texttt{sc\_cheb.m}.
\\

Derivative $d/d\tau$
& \eqref{eq:dt}, \eqref{eq:Dt}
&
\texttt{Dt = D * 2 / (t\_max - t\_min)} in the residual functions.
\\

ODE for $v$
& \eqref{eq:odevt}, \eqref{eq:R1}
&
Residual \texttt{R1} in \texttt{sc\_res\_collapse.m} and
\texttt{sc\_res\_expwave.m}.
\\

ODE for $w$
& \eqref{eq:odewt}, \eqref{eq:R2}
&
Residual \texttt{R2} in \texttt{sc\_res\_collapse.m} and
\texttt{sc\_res\_expwave.m}.
\\

Outer asymptotic boundary condition
& \eqref{eq:bcv}, \eqref{eq:bca}
&
\texttt{sc\_res\_collapse.m}
\\

Critical-point conditions
& \eqref{eq:criticalBC}
&
\texttt{sc\_res\_expwave.m}
\\

Unknown vector $\mathbf u=[\mathbf v;\mathbf w]$
& \eqref{eq:statevector}
&
Nodal values passed to \texttt{sc\_newton\_cs.m}.
\\

Jacobian $J=\partial\mathbf R/\partial\mathbf u$
& \eqref{eq:complexstep}
&
Complex-step differentiation in \texttt{sc\_newton\_cs.m}.
\\

Scaled Newton system
& \eqref{eq:scaled}
&
Row scaling in \texttt{sc\_newton\_cs.m}.
\\

Regularized Newton correction
& \eqref{eq:tikhonov}, \eqref{eq:augmented}
&
Tikhonov-regularized least-squares solve in \texttt{sc\_newton\_cs.m}.
\\

Backtracking
& \eqref{eq:armijo}
&
Damped Newton update in \texttt{sc\_newton\_cs.m}.
\\

Grid refinement in $N$
& \eqref{eq:Ncontinuation}
&
\texttt{sc\_solve\_collapse.m} and \texttt{sc\_solve\_expwave.m}.
\\

Continuation in $A$
& Sec.~\ref{sec:contA}
&
\texttt{sc\_solve\_collapse.m}
\\

Mass relation $m=x^2\alpha(x-v)$
& \eqref{eq:mass}, \eqref{eq:m0numerical}
&
\texttt{shu1977\_spectral\_octave.m}, to compute $m(x)$ and $m_0$.
\\

Comparison with Shu's Table~1
& \eqref{eq:m0numerical}
&
\texttt{shu1977\_spectral\_octave.m}
\\

Comparison with Shu's Table~2
& ---
&
\texttt{shu1977\_spectral\_octave.m}
\\

Validation and error analysis
& Sec.~\ref{sec:validation}
&
Main driver; recorded in the numerical validation output.
\\

\bottomrule
\end{tabular}
\caption{Equation-to-code traceability for the spectral implementation.}
\label{tab:traceability}
\end{table}
This traceability lets a reader follow each mathematical expression
through to the numerical operation that evaluates it.

%=======================================================================
\section{The main driver --- \texttt{shu1977\_spectral\_octave.m}}
%=======================================================================

The main driver \texttt{shu1977\_spectral\_octave.m} connects the solver
components and compares the results with Shu (1977). It performs three
tasks:

\begin{enumerate}

\item \textbf{Comparison with Shu's Table~1 (collapse solutions).}
For
\[
A=2.2,2.4,\ldots,4.0,
\]
the driver calls \texttt{sc\_solve\_collapse}, obtains the corresponding
solution, and evaluates
\begin{equation}
m_0
\approx
m(x_{\min})
=
x_{\min}^{2}\alpha(x_{\min})
\bigl[x_{\min}-v(x_{\min})\bigr].
\label{eq:m0numerical}
\end{equation}
This is an approximation, since $m(x_{\min})\ne m_0$ for finite
$x_{\min}$; the truncation error is estimated by varying $x_{\min}$
(see Sec.~\ref{sec:validation}). The computed values are compared with
those printed in Shu's Table~1.

\item \textbf{Comparison with Shu's Table~2 (expansion-wave solution).}
The expansion-wave calculation uses the critical-point boundary
conditions~\eqref{eq:criticalBC}. The solution is evaluated over
\[
0.001\leq x\leq1
\]
and compared with the values reported in Shu's Table~2, with the
numerical solution interpolated onto the $x$-values of the published
table where necessary.

\item \textbf{Physical reconstruction.}
The driver reconstructs the physical density and velocity fields from
the similarity solution. For the expansion-wave solution, the
self-similar solution is used inside the expansion wave, and the
singular isothermal sphere is used outside it:
\begin{equation}
\rho=\frac{a^2}{2\pi G r^2},
\qquad
u=0.
\end{equation}
The final version of the driver writes the numerical results (Tables~1
and~2, the error analysis, and the validation summary) to the text file
\texttt{shu1977\_results\_spectral\_full.txt}; it does not produce
figures. Figure production is handled separately by the physical
reconstruction utility.

\end{enumerate}

%=======================================================================
\section{Numerical validation}\label{sec:validation}
%=======================================================================

The test is whether the spectral solution reproduces the values
published by Shu (1977). For $m_0$ in Shu's Table~1, the maximum
relative difference is 0.293\% (at $A=2.20$). For Shu's Table~2, the
maximum relative differences are
\[
0.380\% \ \text{for }\alpha,
\qquad
0.618\% \ \text{for }-v,
\qquad
0.442\% \ \text{for }m
\]
at $x=0.20$, $x=0.95$, and $x=0.25$ respectively.
The critical-point values $v(x{=}1)=0$ and $\alpha(x{=}1)=2$ are imposed
as boundary conditions~\eqref{eq:criticalBC}, and $m(x{=}1)=2$ follows
from Eq.~\eqref{eq:mass}; they are therefore not part of the validation. The observed differences can be understood as follows.

\begin{itemize}
\item The maximum relative difference in Table~1, 0.293\% at
      $A=2.20$, and the maxima for $\alpha$ and $m$ in Table~2,
      0.380\% and 0.442\%, are all within the rounding interval of the
      paper's three-significant-figure entries (at most 0.5\%).
\item The maximum relative difference for $-v$, 0.618\%, occurs at
      $x=0.95$, where the paper reports $-v=0.051$ with only two
      significant figures. The rounding interval for this printed
      value is $\pm 0.0005$, i.e.\ a relative interval of up to
      0.98\%. The spectral value is $0.05130$, which lies inside the
      paper's rounding interval. The apparent $0.618\%$ difference is
      therefore entirely due to the paper's printed precision at this
      point.
\item The residual of the Newton iteration converges to $10^{-10}$, so
      the algebraic residual of the spectral solution is not the source
      of the observed differences. The remaining discrepancies are
      dominated by (i) the rounding of the paper's tabulated values and
      (ii) the paper's own numerical accuracy, which is limited by the
      four-step Runge--Kutta integration that Shu uses with a step size
      that may not fully resolve the critical region near $x=1$. Both
      effects are of the same order as the observed differences, so the
      spectral solution may be considered converged relative to the
      precision at which the paper reports its results.
\end{itemize}

%=======================================================================
\section{Summary of the method}
%=======================================================================

The numerical procedure can be summarized as follows:

\begin{enumerate}
\item Rewrite Shu's similarity equations in the variables $\tau=\ln x$
      and $w=\ln\alpha$, and map the $\tau$ interval onto $[-1,1]$.
\item Discretize $v$ and $w$ by their nodal values on the
      Chebyshev--Gauss--Lobatto grid, using the differentiation matrix
      to form the nonlinear residual equations, with the outer-node
      residuals replaced by the boundary conditions.
\item Solve the system with a scaled, Tikhonov-regularized, damped Newton
      iteration whose Jacobian is computed by complex-step
      differentiation.
\item Refine the grid in $N$ and continue in $A$.
\item Recover $\alpha$, $m$ and the physical variables, and compare with
      Shu's Tables~1 and~2.
\end{enumerate}

Table~\ref{tab:methods} contrasts this procedure with that of Shu (1977).

\begin{table}[H]
\centering
\renewcommand{\arraystretch}{1.3}
\begin{tabular}{@{}p{3.4cm}p{5.2cm}p{5.6cm}@{}}
\toprule
 & \textbf{Shu (1977)} & \textbf{This work} \\
\midrule
Formulation & ODEs in $x$, $v$, $\alpha$ & ODEs in $\tau=\ln x$, $v$, $w=\ln\alpha$ \\
Initialization & Asymptotic expansions & Analytic blend of limiting forms \\
Discretization & Four-step Runge--Kutta integration & Chebyshev collocation \\
Nonlinear solver & --- (integration of an initial-value problem) & Complex-step-Jacobian Newton iteration \\
\bottomrule
\end{tabular}
\caption{Methodological differences between Shu (1977) and the present
calculation.}
\label{tab:methods}
\end{table}

%=======================================================================
\section{Source-file summary}
%=======================================================================

\begin{table}[H]
\centering
\renewcommand{\arraystretch}{1.35}
\begin{tabular}{@{}p{5.0cm}p{9.0cm}@{}}
\toprule
\textbf{File} & \textbf{Mathematical / numerical role} \\
\midrule

\texttt{sc\_cheb.m}
&
Constructs the Chebyshev--Gauss--Lobatto nodes and differentiation
matrix $D$.
\\

\texttt{sc\_res\_collapse.m}
&
Constructs the log-transformed collapse residual equations and imposes
the outer asymptotic boundary conditions.
\\

\texttt{sc\_res\_expwave.m}
&
Constructs the log-transformed expansion-wave residual equations and
imposes the critical-point conditions $v(x{=}1)=0$ and
$\alpha(x{=}1)=2$.
\\

\texttt{sc\_newton\_cs.m}
&
Computes the complex-step Jacobian and solves the nonlinear collocation
system using scaled, regularized Newton steps with backtracking.
\\

\texttt{sc\_solve\_collapse.m}
&
Controls the grid refinement in $N$ and the continuation in $A$ for the
collapse family.
\\

\texttt{sc\_solve\_expwave.m}
&
Solves the limiting expansion-wave problem using the critical-point
boundary conditions.
\\

\texttt{shu1977\_spectral\_octave.m}
&
Main driver: comparison with Shu's Tables~1 and~2, validation and error
analysis, and output generation.
\\

\bottomrule
\end{tabular}
\caption{Summary of the source files and their mathematical roles.}
\label{tab:files}
\end{table}

%=======================================================================
\section{Reproducibility statement}
%=======================================================================

The calculation can be repeated with the following information.

\begin{itemize}[nosep]
\item Software: GNU Octave 10.3.0 (or compatible).
\item Code: the main driver \texttt{shu1977\_spectral\_octave.m} together
      with the helper files \texttt{sc\_cheb.m},
      \texttt{sc\_res\_collapse.m}, \texttt{sc\_res\_expwave.m},
      \texttt{sc\_newton\_cs.m}, \texttt{sc\_solve\_collapse.m}, and
      \texttt{sc\_solve\_expwave.m}.
\item Domain: $x_{\min}=10^{-3}$; $x_{\max}=20$ (collapse) and
      $x_{\max}=1$ (expansion wave).
\item Polynomial degrees: $N=30,\,60,\,120$.
\item Collapse parameters: $A=2.2,\,2.4,\ldots,4.0$.
\item Solver parameters: $h=10^{-30}$ (complex step),
      $\epsilon=10^{-10}$ (Tikhonov), $\mathrm{tol}=10^{-10}$,
      $c=10^{-4}$, at most 30 step halvings per iteration.
\item Run time: the full calculation (10 collapse cases plus one
      expansion-wave case) completes in a few seconds on a standard
      desktop machine.
\end{itemize}

\end{document}